\notepage{52}{The reduced radial wavefunction}
Separate the simultaneous eigenfunction as
\[
 \Psi_{E\ell m}(r,\theta,\phi)
 =R_{E\ell}(r)Y_\ell^m(\theta,\phi).
\]
The radial function obeys
\[
 \left\{-\frac{\hbar^2}{2m}
 \left[\frac1{r^2}\frac\dd{\dd r}\left(r^2\frac\dd{\dd r}\right)
       -\frac{\ell(\ell+1)}{r^2}\right]+V(r)\right\}
 R_{E\ell}(r)=E R_{E\ell}(r).
\]
The equation contains no magnetic quantum number $m$, so the radial
function may be chosen independent of that label.

\smallheading{Normalization}
Because the spherical harmonics are normalized,
\begin{align*}
 1&=\int\abs{\Psi_{E\ell m}(\mathbf r)}^2\dd^3r\\
  &=\int_0^\infty r^2\abs{R_{E\ell}(r)}^2\dd r
     \underbrace{\int\abs{Y_\ell^m(\theta,\phi)}^2\dd\Omega}_{1}\\
  &=\int_0^\infty\abs{rR_{E\ell}(r)}^2\dd r.
\end{align*}
Define the reduced radial function
\[
 U_{E\ell}(r)=rR_{E\ell}(r),\qquad
 R_{E\ell}(r)=\frac{U_{E\ell}(r)}r.
\]
Then $\int_0^\infty\abs{U_{E\ell}(r)}^2\dd r=1$, and
\[
 \boxed{\left\{\frac{\dd^2}{\dd r^2}
 +\frac{2m}{\hbar^2}\left[E-V(r)
       -\frac{\ell(\ell+1)\hbar^2}{2mr^2}\right]\right\}
       U_{E\ell}(r)=0.}
\]
The centrifugal term acts as part of the effective one-dimensional
radial potential. Absolute squares above also cover complex radial
functions; ordinary squares are sufficient only when they are chosen
real.
